\documentclass[11pt,a4paper]{article}
\usepackage{fontspec,amsmath,amssymb,amsthm,microtype}
\setmainfont{lmroman10-regular.otf}[BoldFont=lmroman10-bold.otf,ItalicFont=lmroman10-italic.otf,BoldItalicFont=lmroman10-bolditalic.otf]
\usepackage{xcolor}
\usepackage[most]{tcolorbox}
\usepackage{mathrsfs}
\providecommand{\lt}{<}
\providecommand{\gt}{>}
\setmonofont{DejaVuSansMono.ttf}[BoldFont=DejaVuSansMono-Bold.ttf,ItalicFont=DejaVuSansMono-Oblique.ttf,BoldItalicFont=DejaVuSansMono-BoldOblique.ttf,Scale=0.86]
\newfontfamily\leanfallback{STIXTwoMath-Regular.otf}[Scale=0.86]
\definecolor{leanrule}{gray}{0.62}
\definecolor{leanback}{gray}{0.962}
\definecolor{leantitle}{gray}{0.22}
\definecolor{leanmuted}{gray}{0.38}
\newtcolorbox{leanbox}[1]{enhanced,breakable,colback=leanback,colframe=leanrule,
  boxrule=0.4pt,arc=0.6mm,left=2.4mm,right=2.4mm,top=1.8mm,bottom=1.8mm,
  before skip=0.8em,after skip=1em,fonttitle=\small\bfseries,coltitle=leantitle,
  colbacktitle=leanback,titlerule=0.3pt,title={#1}}
\newlength{\leancharwidth}
\newenvironment{leancode}{\par\smallskip\begingroup\ttfamily\footnotesize\raggedright
  \settowidth{\leancharwidth}{M}\setlength{\parskip}{0pt}\setlength{\parindent}{0pt}
  \frenchspacing}{\par\endgroup\smallskip}
\newcommand{\leanline}[2]{\hangindent=\dimexpr#1\leancharwidth+4\leancharwidth\relax
  \hangafter=1\noindent\hspace*{#1\leancharwidth}#2\par}
\newcommand{\leansource}[1]{\par{\footnotesize\color{leanmuted}#1}\par}
\newcommand{\leanlabel}[1]{\par\smallskip\noindent{\small\bfseries #1}\par}
\usepackage[margin=25mm]{geometry}
\usepackage[unicode,hidelinks]{hyperref}
\hypersetup{pdfauthor={LangevinLean Project}}
\setlength{\emergencystretch}{3em}
\setlength{\parskip}{0.4em}
\setlength{\parindent}{0pt}
\allowdisplaybreaks[2]
\raggedbottom
\title{\texorpdfstring{The initial transport infimum}{The initial transport infimum}}
\author{LangevinLean Project}
\date{}
\begin{document}\maketitle
Draft mathematical exposition; these notes may contain mistakes.


The Wasserstein distance is the infimum of quadratic coupling costs. An estimate in terms of an arbitrary initial coupling therefore gives an estimate in terms of the initial Wasserstein distance, provided the comparison process exists for each finite-cost coupling and the numerical output law remains the same. No optimal coupling is required.




Earlier results used here: \cite[Section~3.47]{LangevinLean2026ULASC}.

\section{\texorpdfstring{The infimum over finite-cost initial couplings}{The infimum over finite-cost initial couplings}}


For each finite-cost coupling \(c\) of \(\mu_0\) and \(\pi\), assume an \(L^2\) realization of the initial pair and a common independent Euler noise. Suppose its initial \(L^2\) discrepancy is bounded by \(\operatorname{cost}(c)\) and all realizations have the same numerical output law \(\mu_k\). The synchronous estimate gives \(W_2(\mu_k,\pi)\le q^k\operatorname{cost}(c)+b\) for every \(c\). Taking the infimum gives upper bound \(q^kW_2(\mu_0,\pi)+b\).



Indeed, for each \(\varepsilon>0\) choose an initial coupling of cost at most \(W_2(\mu_0,\pi)+\varepsilon\). The same estimate applies, and the left side and additive term are independent of the choice. Letting \(\varepsilon\) decrease to zero proves the claim. A single finite-cost coupling already proves output finiteness before real-valued transport is used. The Gaussian product construction establishes these assumed realizations and their common numerical marginal.



\begin{thebibliography}{ULA-LSI}
\bibitem[ULA-SC]{LangevinLean2026ULASC}
LangevinLean Project. \emph{Unadjusted Langevin under Strong Convexity}. LL-01, 2026. Draft report, content revision d50c0f64815b. \url{https://odakyildiz.github.io/LangevinLean/docs/pdf/UnadjustedLangevin/FullProof.pdf}.
\end{thebibliography}

\end{document}
